% paper/draft.tex -- SynapForge paper draft skeleton.
%
% Section 4 (Results) is auto-generated by:
%   python scripts/fill_paper_section4.py --log <train.log> \
%       --bench-json bench_results/torch_compile_*.json \
%       --ckpt-step <N> --out paper/section4_auto.tex
%
% Re-run after every rental log scp.
\documentclass[11pt]{article}

\usepackage[margin=1in]{geometry}
\usepackage{booktabs}
\usepackage{graphicx}
\usepackage{hyperref}
\usepackage{amsmath}
\usepackage{amssymb}

\title{SynapForge: A 100M-Parameter Liquid-Spiking Hybrid for Long-Context Inference}
\author{SynapForge Team}
\date{\today}

\begin{document}
\maketitle

\begin{abstract}
SynapForge is a 100M-parameter hybrid recurrent network combining
multiplicative Closed-form Continuous-time (CfC) cells with
Parametric LIF spiking neurons, targeting test-time STDP plasticity
and matmul-free inference. We present training progression, validation
perplexity, and throughput on A800-80GB.
\end{abstract}

\section{Introduction}
\label{sec:intro}
% TODO: motivation, contributions, prior art.

\section{Architecture}
\label{sec:arch}
% TODO: HybridBlock, CfC + PLIF, NeuroMCP, R-fold.

\section{Training}
\label{sec:training}
% TODO: KD setup, data, hyperparameters, phase manager.

% AUTO-GENERATED by scripts/fill_paper_section4.py -- DO NOT EDIT BY HAND.
% Re-run after every train log / bench JSON refresh:
%   python scripts/fill_paper_section4.py --log <log> --bench-json <bench>
%       --ckpt-step <N> --out paper/section4_auto.tex

\section{Results}
\label{sec:results}

\begin{table}[h]
\centering
\caption{Train CE and validation perplexity (TTT vs holdout) at key training steps. \textit{TTT} = test-time-trained val (with inference-time STDP); \textit{holdout} = static val with STDP frozen.}
\label{tab:train_progression}
\begin{tabular}{lrrrr}
\toprule
Step & Train CE & Val PPL (TTT) & Val PPL (holdout) & Notes \\
\midrule
\textbf{TBD} & \textbf{TBD} & \textbf{TBD} & \textbf{TBD} & -- \\
\textbf{TBD} & \textbf{TBD} & \textbf{TBD} & \textbf{TBD} & -- \\
\textbf{TBD} & \textbf{TBD} & \textbf{TBD} & \textbf{TBD} & snapshot \\
\textbf{TBD} & \textbf{TBD} & \textbf{TBD} & \textbf{TBD} & -- \\
\bottomrule
\end{tabular}
\end{table}

\begin{table}[h]
\centering
\caption{Throughput, peak GPU memory, training wall-clock, and \texttt{torch.compile} speedup over the same A800 backbone. Throughput is the mean over all \texttt{step\_ms}-stamped log lines.}
\label{tab:perf}
\begin{tabular}{lr}
\toprule
Metric & Value \\
\midrule
Throughput (tok/s mean) & \textbf{TBD} \\
GPU memory peak (GB) & \textbf{TBD} \\
Training wall-clock (h, est.) & \textbf{TBD} \\
\texttt{torch.compile} speedup & \textbf{TBD} \\
\bottomrule
\end{tabular}
\end{table}

\begin{figure}[h]
\centering
% Place rendered PNG via scripts/plot_train_curves.py here:
% \includegraphics[width=0.95\linewidth]{figures/curves.png}
\caption{Train/val curves (placeholder; rerun \texttt{scripts/fill\_paper\_section4.py} after first rental log scp).}
\label{fig:train_val_curves}
\end{figure}


\section{Discussion}
\label{sec:discussion}
% TODO: limitations, future work.

\bibliographystyle{plain}
% \bibliography{refs}

\end{document}
